\section{Stronger finite certificates}
\label{sec:overnight-finite}
The following refinements strengthen the finite examples without changing
our asymptotic claims. All comparisons in this section use the literal
constants $c_1=c_2=1$: the list prescription is
$2^{H_2(\rho)/\eta}$ and the line prescription is
$n2^{H_2(\rho)/\eta}$. The evaluation domain and alphabet are specified
separately in each example.

\subsection{Exact enumeration at length 64}
\label{sec:overnight-exact64}
An exact enumeration strengthens the finite interval example further.
Among the $34$-subsets $A\subseteq\{0,\ldots,63\}$, exactly
\[
 5{,}552{,}914{,}238{,}035
\]
satisfy $\sum_{a\in A}a=1071$ and
$\sum_{a\in A}\binom a2=22139$.
Their monic root polynomials share the coefficients of degrees $34$, $33$,
and $32$, the last being $\binom{1071}{2}-22139$. Subtracting their common
prefix therefore gives distinct codewords of degree at most $31$, each
agreeing with the prefix word at exactly $34$ coordinates. This proves the
displayed list lower bound over every prime field with $p>64$.
The count uses an exact take-or-omit recurrence in subset size, first moment,
and binomial second moment, with reachability pruning. An independent
ascending-label recurrence on complementary $30$-subsets verifies every
histogram entry; all $129$ previously computed moments also agree.

This is a $5.05\%$ improvement on the degree-40 polynomial-weight
certificate $5285900426578$. At $p=2^{127}-1$, the original strict
below-Elias comparison still applies. The enumeration counts one integer
moment class exactly; it does not claim that this is the largest possible
Reed--Solomon decoding list.

\subsection{Conditional radial weights and a displaced boundary}
\label{sec:overnight-weights}
\begin{proposition}[Conditional signed radial weights]
Let $\mathcal A$ be a family of $t$-subsets of
$\{0,\ldots,n-1\}$ sharing their first binomial moment, and put
$C=|\mathcal A|$. For $j=2,3,4,5$, choose
\[
 b_j(x)=\binom{x}{j}+\sum_{i<j}c_{ji}\binom{x}{i},
 \qquad c_{ji}\in\mathbb Q.
\]
For a uniform $A\in\mathcal A$, write
$B_j=\sum_{a\in A}b_j(a)$. Let $U_2,\ldots,U_5$ be independent
uniform variables on $[-1/2,1/2)$, independent of $A$, and choose
positive rational scales $s_j$. Set
\[
 X_j=\frac{B_j-\mathbb EB_j+U_j}{s_j},\qquad
 S=\sum_{j=2}^5X_j^2,\qquad \mu_r=\mathbb ES^r.
\]
For $w\in\mathbb Q[T]$ with integrable positive part on $[0,\infty)$,
suppose $E=\mathbb E[w(S)]>0$. Then some class sharing all five
binomial moments has size at least
\[
 \left\lceil
 \frac{CE}{\pi^2(\prod_{j=2}^5s_j)
       \int_0^\infty w(u)_+u\,du}
 \right\rceil.
\]
Consequently, for every prime $p>n$, the ordinary Reed--Solomon code
on $\{0,\ldots,n-1\}$ with dimension $t-5$ has a list of that size
at agreement threshold $t$.
\end{proposition}

\begin{proof}
Because the size and first moment are fixed, the remaining moment vectors
lie in an affine lower-triangular lattice with diagonal entries one.
Its translates of $[-1/2,1/2)^4$ tile $\mathbb R^4$: successively
rounding coordinates determines the lattice point uniquely. If $L$ is
the largest fiber, the density of $B+U$ is at most $L/C$, and that of
$X$ is at most $(L/C)\prod_js_j$.

Radial integration in four dimensions gives
\[
 \int_{\mathbb R^4}
  w(\|x\|^2)_+\,dx
 =\pi^2\int_0^\infty w(u)_+u\,du.
\]
Replacing the signed weight by its positive part and applying the density
bound proves the inequality. Triangular inversion fixes the first five
integer binomial moments. Newton identities then fix the leading six
terms of the monic support locators. Subtracting those locators from
their common prefix gives distinct polynomials of degree less than
$t-5$, agreeing exactly on their supports. Distinctness after reduction
follows from uniqueness of the monic locator of a subset of distinct
field elements.
\end{proof}

\paragraph{Choosing the weight.}
The subclass $w(u)=(R-u)P(u)^2$ is convenient. For
$P(u)=\sum_{i=0}^dc_iu^i$, its expectation and radial integral are
quadratic forms with matrices
\[
 A_{ij}(R)=R\mu_{i+j}-\mu_{i+j+1},\qquad
 B_{ij}(R)=\frac{R^{i+j+3}}{(i+j+2)(i+j+3)}.
\]
The matrix $B(R)$ is positive definite. At fixed $R$, its generalized
Rayleigh quotient gives the largest ratio in this polynomial class.
Numerical eigenvectors propose a weight; rational rounding and exact
evaluation certify the selected weight independently of the optimization.

\paragraph{A verified M31 instance.}
Take $p=2^{31}-1$, $n=157$, $t=68$, and condition on
$\sum_{a\in A}a=5304$. The conditional population is
\[
 C=4037124019196117766480536278705909334612974.
\]
Using normalized population Gram polynomials of degrees $2$ through $5$,
the scales
\[
 (s_2,s_3,s_4,s_5)=(5751,76272,754066,5946077),
\]
and a degree-seven weight
\[
 w(u)=-\prod_{i=1}^7(u-r_i),
\]
with positive ordered rational roots approximately
\[
 (r_1,\ldots,r_7)\approx
 (2.78693,5.20044,5.74234,11.14408,11.19106,20.21394,20.21451),
\]
gives a list of at least
\[
 28165503856884755099.
\]
Exact rational logarithm intervals verify that the radius is strictly
below Elias and that the list exceeds
$2^{H_2(63/157)/(5/157)}$ by more than $2^{34.09946}$.
The positive intervals are $[0,r_1]$, $[r_2,r_3]$, $[r_4,r_5]$, and $[r_6,r_7]$,
so their integrals are exact polynomial endpoint differences.
The certificate uses $280$ scalar directions through degree fourteen and
exact cube-noise moments. Of $560$ subset-mode outputs, $351$ were counted
directly and $209$ obtained through independently checked reflections.
All complement, prior-moment, and projection-polynomial identities pass. The exact rational roots, rather than their displayed decimal
approximations, are in \path{research/overnight_2026-09-16/certificates/radial_moment_optimum7_verification.json}.

\paragraph{Limit of the truncated radial calculation.}
An exact rational contraction certificate proves existence of ordered
roots and a constant $\lambda>0$ satisfying
\[
 \mu_j=\frac{\lambda}{j+2}
       \sum_{i=1}^7(-1)^{i+1}r_i^{j+2},\qquad 0\le j\le7.
\]
The density equal to $\lambda/\pi^2$ on these four radial bands has
the prescribed moments. It therefore supplies an upper limit on every
degree-at-most-seven radial weight certificate; the corresponding signed
root polynomial attains it. The validated interval identifies the same
rounded integer bound displayed above. This is an optimum of the
truncated radial moment method, not an upper bound on the actual RS list.

\begin{proposition}[Displacing a radial boundary]
Let $X\in\mathbb R^d$ have density at most $M$, where $d\in\{3,4\}$.
Let $0<r_1<\cdots<r_7$, put
$l(s)=\prod_{j=2}^7(s-r_j)$, and define
\[
 w_0(s)=l(s)(r_1-s),\qquad
 w_v(x)=l(|x|^2)(r_1-|x|^2+v\cdot x).
\]
Suppose $|v|^2r_2<(r_2-r_1)^2$. Write $a=v/2$ and
$R=r_1+|a|^2$. Then
\begin{align*}
 \int_{\mathbb R^d}(w_v)_+
 &=\int_{\mathbb R^d}(w_0(|x|^2))_+
   -\int_{|x|^2<r_1}w_0(|x|^2)\,dx\\
 &\quad+\int_{|y|^2<R}(R-|y|^2)l(|y+a|^2)\,dy.
\end{align*}
In particular, if $\mathbb E w_v(X)>0$, then
$M\ge\mathbb E w_v(X)/\int(w_v)_+$.
\end{proposition}

\begin{proof}
Completing the square gives
$r_1-|x|^2+v\cdot x=R-|x-a|^2$.
The assumed inequality places this displaced ball strictly inside
$\{|x|^2<r_2\}$. The other positive annuli remain unchanged. On each
origin-centered annulus the added term $(v\cdot x)l(|x|^2)$ integrates
to zero. Substituting $x=y+a$ on the displaced ball proves the integral
identity. Finally, $\mathbb E w_v\le\mathbb E(w_v)_+\le M\int(w_v)_+$.
\end{proof}

For the normalized smoothed moment lattice, $M\le(L/C)\prod_js_j$.
The required expectations use no additional subset counts: differentiating
an identity $|x|^{2j+2}=\sum_u c_u(u\cdot x)^{2j+2}$ gives
$x_i|x|^{2j}=\sum_u c_u u_i(u\cdot x)^{2j+1}$.
The displaced-ball integral is exact by angular moments or Cartesian
monomial integration. At $n=157$, $k=63$, $t=68$, the verified
four-dimensional certificate gives $L\ge28169451256663519418$ over
$\mathbb F_{2^{31}-1}$, with finite-prescription excess greater than
$34.09967$ bits. The three-dimensional certificate gives source class
$508108952862977950462384886$ and quartic-extension line bank
$508108952659354448051620618$, with excess greater than $50.90955$ bits.
The selected rational vectors, exact integrals, and separate raw-moment
checks are recorded in the ancillary displaced-boundary certificates.

For the three-dimensional version, retain Gram coordinates of degrees
$2,3,4$ after conditioning on the first moment. The radial integration
factor is $2\pi$, and its radial density is $s^{1/2}$. The same triangular
lattice argument certifies a class with four fixed integer moments.
Its locator prefix is fixed through degree $64$, so its residuals have
degree at most $63$. The rational-line lemma lowers that degree to $62$.
The quartic-extension instance therefore has length $157$, dimension $63$,
and $68$ agreements, with common agreement exactly $63$ and selected-graph
concurrency at most $18$. Its evaluation points remain in $\mathbb F_p$;
its received words and labels use $\mathbb F_{p^4}$. The Elias comparison
uses the characteristic $p=2^{31}-1$.

\subsection{Exact incidences and prime-alphabet lines}
\label{sec:overnight-prime-lines}
\begin{proposition}[Coordinate and next-coefficient collision savings]
Let $N$ distinct support locators share a prefix leaving residual degree
at most $k$ over $\mathbb F_p$, on an evaluation domain of size $n<p$.
Let $I_x$ count the supports containing $x$. If their next coefficients
take at most $K$ values, then some pole outside the domain has at most
\[
 \left\lfloor\frac{
 k\binom N2-\sum_x\binom{I_x}{2}-P(N,K)}{p-n}\right\rfloor
\]
colliding unordered pairs, where
$P(N,K)=K\binom a2+ab$ for $N=aK+b$, $0\le b<K$.
\end{proposition}

\begin{proof}
At least $P(N,K)$ pairs share the next coefficient and have difference
degree at most $k-1$; the others have degree at most $k$.
Shared support roots account for exactly $\sum_x\binom{I_x}{2}$ roots
of these differences on the evaluation domain. Subtracting them and
averaging over the remaining $p-n$ poles proves the assertion.
\end{proof}

The least $J$ for which $P(N,J)$ does not exceed this pair budget is a
guaranteed number of distinct line labels. For the interval $\{0,\ldots,81\}$,
the selected twelve-subset class has first sum $486$, second binomial
moment $12792$, and $N=221190005$. Its exact coordinate incidences and
a rationally certified next-moment range give $J=138752510$ over
$\mathbb F_{2^{31}-1}$, with $k=9$, $t=12$, common agreement $9$,
and concurrency at most $24$. This exceeds the stated finite line
prescription by more than $2^{7.04613}$. The largest two-moment class
has more supports, but yields a weaker line certificate.

Here the source locators share the prefix
$X^{12}-486X^{11}+105063X^{10}$. The exact shared-root count is
$\sum_x\binom{I_x}{2}=43128963788215318$. Write
$Z=\sum_{x\in A}\binom{x}{3}$. Newton identities give
\[
 e_3=\binom{486}{3}+(2-486)12792+2Z.
\]
Exact rational dual inequalities place $Z$ between $222145$ and $276503$.
Explicitly, for $\sigma\in\{1,-1\}$ and rational multipliers
$\lambda=(\lambda_0,\lambda_1,\lambda_2)$, every selected support satisfies
\[
 \sigma Z\ge 12\lambda_0+486\lambda_1+12792\lambda_2
 +\sum_{x=0}^{81}\min\left\{0,
 \sigma\binom{x}{3}-\lambda_0-\lambda_1x-\lambda_2\binom{x}{2}\right\}.
\]
This follows termwise from $z_xr_x\ge\min\{0,r_x\}$ for
$z_x\in\{0,1\}$. The choices
$(\sigma,\lambda)=(1,(1421,-509,107/3))$ and
$(-1,(-9568,2447/3,-130/3))$ give the exact lower bounds
$666433/3$ and $-829511/3$, respectively. Integer rounding gives
the stated interval for $Z$.
Thus at least $449907051360$ source pairs share the next coefficient,
and the averaged allowed-pole collision budget is $82437495$.
Inverting the balanced-pair function gives the stated $138752510$ labels.
The coordinate counts have a separate formal-division check, using the
same underlying histogram engine in a different order. The next-coefficient
bounds use rational termwise inequalities, independently of the exploratory
minimum/maximum recurrence.

\subsection{Stronger native circle examples}
\label{sec:overnight-circle}
\begin{proposition}[A refined circle-pencil bank]
Let $N$ distinct residual polynomials $R_A$ be obtained from the
anchored $h$-subsets of $G=\mu_M$ in one reciprocal coefficient class,
and put $d=h-2r-2$. For integers $v\ge0$ and $b\ge1$, define
\[
 \Pi(v,b)=(b-e)\binom a2+e\binom{a+1}2,
 \qquad v=ab+e,\quad 0\le e<b.
\]
Set
\[
 C=d\binom N2-\binom N2-\Pi(N(h-1),M-1)
       -(p+1-M)\Pi(N,p).
\]
For poles in $\mathbb F_{p^2}\setminus\mu_{p+1}$ take
$Q=p^2-p-1$; for poles in
$\mathbb F_{p^4}\setminus\mathbb F_{p^2}$ take $Q=p^4-p^2$.
There is an admissible pole with at least
\[
 \min\{j\in\{1,\ldots,N\}:\Pi(N,j)\le\lfloor C/Q\rfloor\}
\]
distinct labels. The agreement and absence-of-correlated-agreement
conclusions of the original construction are unchanged.
\end{proposition}
\begin{proof}
The function $\Pi(v,b)$ is the minimum number of colliding unordered
pairs in a partition into at most $b$ classes. For each $y\in G$, let
$s_y$ count supports containing $y$. Every shared support point is
a root of $R_A-R_{A'}$, so double counting gives
\[
 \sum_{A<A'}|A\cap A'|
 =\sum_{y\in G}\binom{s_y}2
 \ge\binom N2+\Pi(N(h-1),M-1).
\]
At every norm-one point, reciprocity gives
$R_A(y)^p=c^{-1}y^{-d}R_A(y)$ with the same $c$ for the entire class.
Thus at most $p$ evaluation values occur there. Each of the other
$p+1-M$ norm-one points contributes at least $\Pi(N,p)$ pair roots.
These points and $G$ are disjoint. Subtracting their contributions
from the total degree budget $d\binom N2$ bounds the sum of pair
collisions over admissible poles by $C$. Some pole has at most
$\lfloor C/Q\rfloor$ collisions. Balancing its label multiplicities
proves the assertion.
\end{proof}

For $p=2^{31}-1$, $M=128$, $h=19$, and $r=0$, this gives
\[
 J\ge365153906390146587
\]
with quadratic-field poles, improving the earlier
$285003988493147037$. At length $2^{26}$ and rate $1/8$, the bank
exceeds the compared finite prescription by more than $2^9$.
At length $2048$, dimension $256$, and $302$ agreements, it exceeds
that prescription by more than $2^{23}$. Both radii are strictly
below the characteristic-based Elias bound. These numerical
statements have exact integer and rational certificates in
\path{research/overnight_2026-09-16/certificates/circle_threshold_verification.json}.

\paragraph{Counting omitted-support collisions as well.}
There is a further exact saving at each $y\in G$. Exactly $s_y$ residuals
take the value $v(y)$, whereas the others occupy at most $p-1$ values
in the same one-dimensional prime-field subspace. Their collision count
is therefore at least
\[
 f_N(s_y)=\binom{s_y}{2}+\Pi(N-s_y,p-1).
\]
This counts disjoint pairs and can replace the support-root term above.
Its forward difference is
$s-\lfloor(N-s-1)/(p-1)\rfloor$, an increasing function, so balancing
unknown incidences gives a valid lower bound. At reserve $r=0$, an exact
cyclic subset-product count supplies the class population and every $s_y$.
For $M$ dyadic, if $C_k(c)$ counts $k$-subsets of $\{1,\ldots,M-1\}$
with exponent sum $c\pmod M$, then
\[
 C_k(c)=\frac1M\sum_{a\mid M}\mathcal R_a(c)
 (-1)^{k+\lfloor k/a\rfloor}\binom{M/a-1}{\lfloor k/a\rfloor},
\]
where $\mathcal R_a$ is the Ramanujan sum. This follows by extracting the
coefficient of $z^k$ in $(1-(-z)^a)^{M/a}/(1+z)$ for each character.
For a nonanchor exponent $b$, removing its generating factor gives
\[
 I_b=\sum_{j=0}^{h-2}(-1)^j C_{h-2-j}(c-(j+1)b).
\]
Applying these exact populations and collisions gives the following banks.
All rows have $p=2^{31}-1$, native circle domains, and radii strictly below
the characteristic-based Elias bound.
\begin{center}
\begin{tabular}{crrrrl}
\toprule
Alphabet & $n$ & $K$ & $T$ & Guaranteed labels & Excess bits\\
\midrule
$\mathbb F_{p^2}$ &2048&256&302&365153907657996934&$>23$\\
$\mathbb F_{p^2}$ &4096&64&86&334185312123996677&$>24.5949$\\
$\mathbb F_{p^4}$ &8192&1024&1182&73497635817388704528460&$>34.77733$\\
\bottomrule
\end{tabular}
\end{center}
The first row improves the same-parameter quadratic bank in
Appendix~\ref{app:compact-certificates} by about $28.1\%$.
The last two rows maximize the checked guarantee within the stated
one-anchor dyadic family, using complete product classes at $r=0$ and
generic prefix pigeonholing at $r>0$. This is an optimization of that
formula, not of arbitrary circle lists or actual coefficient classes.
The checked range is dyadic $M\ge4$, $B\ge2$, $MB\le2^{30}$,
dyadic packet dimensions $K/B<M$, and every reserve satisfying the
strict Elias condition and the construction's support-size requirements.
A separate coefficient-sum implementation checks all $703$ computed
product classes and $1406$ collision certificates.

\paragraph{Scope and reproducibility.}
The interval-domain results and the native circle results above are
different constructions. None lowers the pinned better.codes display of
$116.13$, and none is a revised security bound for the complete deployed
Stwo protocol. Their contribution is a stronger set of counterexamples
to the specified finite coding prescriptions. The proof statements,
portable inputs, and verification commands are in
\path{research/overnight_2026-09-16/README.md}.
